\documentclass[11pt,a4paper]{article}
\usepackage[margin=25mm]{geometry}
\usepackage{amsmath,amssymb,amsthm}
\usepackage[T1]{fontenc}
\usepackage{lmodern}
\usepackage[colorlinks=true,linkcolor=blue,citecolor=blue,urlcolor=blue]{hyperref}
\newtheorem{theorem}{Theorem}
\newtheorem{lemma}[theorem]{Lemma}
\newtheorem{proposition}[theorem]{Proposition}
\newcommand{\R}{\mathbb R}
\newcommand{\N}{\mathbb N}
\newcommand{\Tv}{T_v}
\setlength{\emergencystretch}{1em}
\title{A vertex maximum principle for subcritical reinforcement}
\author{Alec Kriebel\\\small Independent researcher\\
\small ORCID: \href{https://orcid.org/0009-0001-9320-500X}{0009-0001-9320-500X}}
\date{13 October 2026\quad---\quad Version 1.0}
\begin{document}
\maketitle
\begin{abstract}
We give a self-contained alternative proof of positive-equilibrium uniqueness
and almost-sure coordinatewise homogenization for subcritical reinforcement on
countable graphs of bounded degree. Positive vertex firing rates are bounded
above but may have infimum zero. The proof combines an exact compensated jump
transform with a maximum principle for bounded positive complete trajectories
in vertex coordinates. Initial edge counts are arbitrary finite positive
integers. Qin has also proved the full homogenization theorem, with broader
positive real initialization; that result is prominently credited here.
Unrestricted uniqueness among nonnegative equilibria is false, and is kept
separate from the corrected positive statement. No first-priority claim is made.
\end{abstract}

\section{The question and the precise result}

In the graph-based reinforcement model, each vertex has an independent
Poisson clock. At a ring, one incident edge is chosen proportionally to its
current weight raised to a power $\alpha$, and its weight increases by one.
Couzini\'e and Hirsch~\cite{CH} proved homogenization on general bounded-degree
graphs for $\alpha<1/2$, slightly beyond $1/2$ on regular graphs, and throughout
the subcritical range on $\mathbb Z$. Earlier heterogeneous-rate results and
the all-subcritical conjecture appear in Couzini\'e's thesis~\cite{thesis}.
The question for bounded positive heterogeneous firing rates is stated in the
Oberwolfach report~\cite{OWR} and the official pre-publication version of the
MATRIX chapter~\cite{MATRIX}.

The report's literal equilibrium definition does not require positive
coordinates. Unrestricted uniqueness is false. On $\mathbb Z$, with unit
vertex rates and $0<\alpha<1$, put weight two on a perfect matching and zero
on all other edges. Every incident denominator is $2^\alpha$. Each positive
edge receives rate two, and each zero edge receives rate zero. Both phases
and the constant-one array are equilibria. This is the same local matching
mechanism as the older normalized alternating-cycle allocation in finite
WARM~\cite[\S3.6]{WARM}; its validity for every positive exponent follows by
direct substitution. Qin~\cite[Remark 1.3]{Qin} also records the infinite
matching example. We claim no novelty for this boundary construction.

The positive theorem below is an explicitly corrected homogenization
formulation, not a proof of that false unrestricted uniqueness sentence.
Qin~\cite[Theorem 1.1]{Qin} proves this theorem and permits arbitrary finite
positive real initial weights. His connected formulation extends componentwise
to disconnected countable graphs. Our initialization class is smaller.
The contribution here is a different proof: finite-volume vertex minimization,
an exact jump compensator, and an approximate space-time maximum principle.
The power/incidence reduction and complete-limit strategy are shared structural
ideas, and are not asserted to originate here. Qin also reports independent
Angel--Hermon--Yao work; its full manuscript and public date were not verified
in our bounded literature audit.

Let $G=(V,E)$ be a countable undirected simple graph. Discard isolated vertices
and assume
\[
 1\leq d(v)\leq D<\infty,\qquad 0<p_v\leq P<\infty.
\]
There is no positive lower bound on all $p_v$. For an edge array $w>0$, set
\[
 S_v(w)=\sum_{e\ni v}w_e^\alpha.
\]
The clocks have rates $p_v$, and choose $e\ni v$ with probability
$N_e^\alpha/S_v(N)$. Initially $N_e(0)=n_e\in\N$, with $n_e\geq1$ finite for
each edge; no uniform bound on $(n_e)$ is imposed. Write $X_e(t)=N_e(t)/t$.

\begin{theorem}\label{thm:main}
For $0\leq\alpha<1$, there is exactly one equilibrium positive on every edge:
\begin{equation}\label{eq:edge}
 x_{uv}=x_{uv}^{\alpha}
 \left(\frac{p_u}{S_u(x)}+\frac{p_v}{S_v(x)}\right).
\end{equation}
Almost surely, $X_e(t)\longrightarrow x_e$ simultaneously for all $e\in E$.
The convergence is coordinatewise.
\end{theorem}

The empty edge set gives the unique empty array and vacuous convergence.
For the proof assume $E\ne\varnothing$ and $0<\alpha<1$, and put
$c=1-\alpha$ and $\beta=\alpha/c>0$. The endpoint $\alpha=0$ is treated at the
end. No uniform spatial convergence, stochastic rate, zero-initial-edge,
critical-exponent or unbounded-degree extension is claimed.

\section{Positive equilibria in vertex coordinates}

For a positive equilibrium define $z_v=p_v/S_v(x)$. Dividing
\eqref{eq:edge} by $x_{uv}^{\alpha}$ gives the bijective correspondence
\begin{equation}\label{eq:vertex}
 x_{uv}=(z_u+z_v)^{1/c},\qquad
 z_v\Tv(z)=p_v,\qquad
 \Tv(z)=\sum_{w\sim v}(z_v+z_w)^\beta.
\end{equation}
Every positive vertex solution obeys
\begin{equation}\label{eq:zupper}
 z_v\leq(p_v/d(v))^c\leq B:=P^c,
\end{equation}
since $\Tv(z)\geq d(v)z_v^\beta$. Thus boundedness is a deduction for every
positive equilibrium, rather than an extra hypothesis.

\begin{proposition}\label{prop:equilibrium}
The vertex equations in \eqref{eq:vertex} have exactly one positive solution.
\end{proposition}
\begin{proof}
For a finite vertex set $F$, set $z_w=0$ outside $F$ and minimize
\[
 \Phi_F(z)=\sum_{\{u,v\}\in E:\{u,v\}\cap F\ne\varnothing}
 \frac{(z_u+z_v)^{\beta+1}}{\beta+1}
 -\sum_{v\in F}p_v\log z_v
 \qquad(z_v>0, v\in F).
\]
Only finitely many edges occur. The power sum dominates a positive multiple of
$\sum_{v\in F}z_v^{\beta+1}$ and hence the logarithms at infinity. On bounded
sublevel sets, approaching a coordinate hyperplane costs infinity. An interior
minimum exists. The logarithmic terms are strictly convex and the power terms
are convex, so the minimum is unique. Its equations are precisely
$z_v\Tv(z)=p_v$ for $v\in F$.

The upper estimate \eqref{eq:zupper} holds, as does
\[
 z_v\geq \frac{p_v}{d(v)(2B)^\beta}>0\qquad(v\in F).
\]
Exhaust $V$ by finite sets and take a diagonal subsequence of their minimizers.
Each fixed coordinate eventually has this positive lower bound and the common
upper bound $B$. Every vertex equation is a finite sum, so it passes to the
coordinatewise limit. This proves positive infinite-volume existence.

For uniqueness let $z,y$ be positive solutions and suppose
$M=\sup_v|z_v-y_v|>0$. Interchange them if necessary and choose, for arbitrarily
small $\varepsilon>0$, a vertex with $z_v-y_v>M-\varepsilon$. Each adjacent
pair sum for $z$ is at least the corresponding sum for $y$ minus $\varepsilon$.
Let $\omega$ be a modulus of continuity of $r^\beta$ on $[0,2B]$. Then
$\Tv(z)-\Tv(y)\geq-D\omega(\varepsilon)$, while
$z_v\geq M-\varepsilon$ and $\Tv(z)\geq(M-\varepsilon)^\beta$. Consequently
\[
 0=(z_v-y_v)\Tv(z)+y_v[\Tv(z)-\Tv(y)]
 \geq(M-\varepsilon)^{\beta+1}-BD\omega(\varepsilon),
\]
which is positive for small $\varepsilon$. No maximum attainment or uniform
positive lower bound is used.
\end{proof}

\section{An exact stochastic transform}

The graphical construction is local. A fixed query up to time $T$ has causal
predecessors along vertex-clock paths with at most $D+1$ choices at each step
and strictly decreasing times. The expected number of length-$m$ ordered paths
is bounded by a fixed local factor times $((D+1)PT)^m/m!$, tending to zero.
The predecessor tree is finitely branching; infinite ancestry would give an
infinite path. Thus local queries have finite ancestry almost surely. Marks
resolve those finite histories. Countably many edge queries and integer
horizons give a consistent process, even when the total firing rate is infinite.
This clock construction does not require bounded initial counts.

For $e=\{u,v\}$, endpoint-clock strong laws imply
\begin{equation}\label{eq:uppergrowth}
 \limsup_{t\to\infty}X_e(t)\leq p_u+p_v\leq2P.
\end{equation}
Use predictable left limits to define
\[
 q_e(t)=\frac{p_u}{S_u(N(t-))}+\frac{p_v}{S_v(N(t-))},\qquad
 Q_e(t)=\int_0^t q_e(s)\,ds.
\]
The edge count has intensity $N_e(t-)^\alpha q_e(t)\leq p_u+p_v$.
For positive integers let
\[
 H(n)=\sum_{j=1}^{n-1}j^{-\alpha}=\frac{n^c}{c}+O(1).
\]
Its exact jump is $n^{-\alpha}$. Therefore
\begin{equation}\label{eq:compensator}
 H(N_e(t))-H(n_e)=Q_e(t)+M_e(t),\qquad
 \langle M_e\rangle(t)=\int_0^t N_e(s-)^{-\alpha}\,dQ_e(s)\leq Q_e(t),
\end{equation}
where $M_e(0)=0$. Since $N_e\geq1$ and $q_e\leq2P$, all finite-horizon
martingales here are square integrable.

Whenever $Q_e\to\infty$, \eqref{eq:compensator} gives $M_e/Q_e\to0$ almost
surely. Indeed, with $A=1+Q_e$, the martingale
$L(t)=\int_0^t A(s)^{-1}\,dM_e(s)$ has bracket at most
$\int_0^\infty(1+Q_e)^{-2}\,dQ_e\leq1$. It is bounded in $L^2$ and converges
almost surely. Integration by parts gives
\[
 \frac{M_e(t)}{A(t)}=L(t)-\frac{1}{A(t)}\int_0^t L(s)\,dA(s).
\]
On $A(t)\to\infty$, weighted Cesaro convergence makes the two terms have the
same limit. This argument does not assume independence of $Q_e$ and $M_e$.

Fix $\varepsilon>0$. By \eqref{eq:uppergrowth}, eventually all finitely many
edges incident to $u$ or $v$ have counts at most $(2P+\varepsilon)t$. Thus
\[
 q_e(t)\geq k_{e,\varepsilon}t^{-\alpha},\qquad
 k_{e,\varepsilon}=\frac{p_u/d(u)+p_v/d(v)}{(2P+\varepsilon)^\alpha}>0.
\]
This first proves $Q_e\to\infty$ and
$\liminf Q_e(t)/t^c\geq k_{e,\varepsilon}/c$. The transform and its normalized
martingale law then imply $N_e(t)^c/(cQ_e(t))\to1$, and hence
\begin{equation}\label{eq:lowergrowth}
 \liminf_{t\to\infty}X_e(t)\geq
 m_e:=\left(\frac{p_u/d(u)+p_v/d(v)}{(2P)^\alpha}\right)^{1/c}>0.
\end{equation}
The upper bound also gives $H(N_e(t))=O(t^c)$, so $Q_e(t)=O(t^c)$ and
$M_e(t)/t^c\to0$. All statements hold on one probability-one event by
countability. The constants and random starting times need not be uniform in
space. The fixed initial term $H(n_e)$ vanishes after scaling.

\section{Complete limiting trajectories}

Define nonnegative vertex primitives
\[
 A_v(t)=ct^{-c}\int_0^t\frac{p_v}{S_v(N(s))}\,ds.
\]
Left limits are immaterial in this Lebesgue integral. Splitting $Q_e$ in
\eqref{eq:compensator} gives
\begin{equation}\label{eq:representation}
 X_{uv}(t)^c=A_u(t)+A_v(t)+r_{uv}(t),\qquad r_{uv}(t)\longrightarrow0.
\end{equation}
Nonnegativity and \eqref{eq:uppergrowth} imply
$\limsup A_v(t)\leq K:=(2P)^c$. Integrating the local incident-count upper
bound similarly gives
\[
 \liminf_{t\to\infty}A_v(t)\geq
 \ell_v:=\frac{p_v}{d(v)(2P)^\alpha}>0.
\]
In logarithmic time $a_v(s)=A_v(e^s)$, differentiation holds almost everywhere:
\begin{equation}\label{eq:prelimit}
 a_v'(s)=c\left[\frac{p_v}{S_v(X(e^s))}-a_v(s)\right].
\end{equation}
The finitely many incident-edge lower bounds \eqref{eq:lowergrowth} bound this
derivative for each fixed vertex at all sufficiently large times. Thus the
translated coordinate functions are eventually uniformly Lipschitz on each
compact interval, with a bound allowed to depend on the coordinate.

For any $s_n\to\infty$, Arzela--Ascoli and a diagonal subsequence over the
countably many vertices and intervals $[-j,j]$ give locally uniform coordinate
limits $a_v(s_n+s)\to b_v(s)$, defined for every $s\in\R$, with
\[
 \ell_v\leq b_v(s)\leq K.
\]
For each edge the error in \eqref{eq:representation} is uniformly small on a
shifted compact logarithmic interval, because its minimum physical time tends
to infinity. Since $X_{uv}^\alpha=(X_{uv}^c)^\beta$, finite incident sums and
positive limiting coordinates justify passing through the denominators in the
integral form of \eqref{eq:prelimit}. Every limit is a classical complete
trajectory of
\begin{equation}\label{eq:limit}
 b_v'=c\left[R_v(b)-b_v\right],\qquad
 R_v(b)=\frac{p_v}{\Tv(b)}.
\end{equation}
This is product-topology compactness. No spatially uniform random time,
supremum-norm compactness or infinite-dimensional stochastic approximation
theorem has been assumed.

\section{The vertex maximum principle and convergence}

\begin{lemma}\label{lem:rigidity}
Every positive complete classical solution of \eqref{eq:limit} with a common
finite upper bound is the equilibrium $z$ of Proposition~\ref{prop:equilibrium}.
No common positive lower bound is required.
\end{lemma}
\begin{proof}
Choose a common upper bound $K_0\geq B$ for $b$ and $z$, and suppose
$M=\sup_{v,s}|b_v(s)-z_v|>0$. Set $d_v=b_v-z_v$. Let $\omega$ be a modulus
for $r^\beta$ on $[0,2K_0]$. Choose $0<\varepsilon<M/4$ so that
\[
 \frac{BD\omega(\varepsilon)}{(M/2)^\beta}<M/4.
\]
If $d_v\geq M-\varepsilon$, every neighbor has $d_w\geq-M$, whence
$\Tv(b)\geq\Tv(z)-D\omega(\varepsilon)$. Also $b_v>M/2$, so
$\Tv(b)\geq(M/2)^\beta$. Using $p_v=z_v\Tv(z)$ gives
\[
 R_v(b)-z_v=\frac{z_v[\Tv(z)-\Tv(b)]}{\Tv(b)}<M/4,
 \qquad d_v'\leq-cM/2.
\]
If $d_v\leq-M+\varepsilon$, every neighbor has $d_w\leq M$, whence
$\Tv(b)\leq\Tv(z)+D\omega(\varepsilon)$. Here $z_v>M/2$, so
$\Tv(z)\geq(M/2)^\beta$. Reciprocal monotonicity gives
\[
 R_v(b)-z_v\geq
 -\frac{z_vD\omega(\varepsilon)}{\Tv(z)+D\omega(\varepsilon)}>-M/4,
 \qquad d_v'\geq cM/2.
\]
In particular, the lower-strip estimate needs no uniform lower bound on the
current denominator $\Tv(b)$.

Some fixed coordinate at some time lies strictly in one extreme strip.
In the upper case, take the connected component of
$\{s:d_v(s)>M-\varepsilon\}$ containing that time. A finite left endpoint
would have value $M-\varepsilon$; integrating the negative derivative from it
contradicts entry into the strip. The component therefore extends to minus
infinity. Integrating backward gives arbitrarily large $d_v$, contradicting
$d_v\leq M$. The lower case is identical with the sign reversed. Thus $M=0$.
Uniform continuity suffices even for $0<\beta<1$; neither an attained maximum
nor differentiation of an infinite supremum is used.
\end{proof}

Every limiting trajectory from the previous section satisfies the lemma and
equals $z$. Hence $a_v(s)\to z_v$ for each vertex: a sequence separated from
$z_v$ would have a limiting trajectory separated at time zero, a contradiction.
Equation~\eqref{eq:representation} then gives
$X_{uv}(t)\to(z_u+z_v)^{1/c}$. Proposition~\ref{prop:equilibrium} identifies
this as the unique positive equilibrium. For $\alpha=0$, independent marking
and Poisson thinning directly give
\[
 X_{uv}(t)\longrightarrow p_u/d(u)+p_v/d(v),
\]
which is also the unique positive equilibrium. This completes the proof of
Theorem~\ref{thm:main}.

\section*{Chronology, verification and disclosure}

The candidate argument is recorded in repository PR~764~\cite{candidate},
whose current GitHub server creation and update timestamps are 5 October 2026;
Qin's arXiv submission is dated 8 October 2026. These observations do not
certify when every current proof byte first became publicly available, or
historical priority over all work. The available literature now contains the
full positive theorem. Our note presents an alternative proof and makes no
first, sole, new-theorem or currently-open claim. The AHY manuscript reported
by Qin remains a bibliographic limitation, not a purportedly inspected proof.
The final publisher MATRIX chapter body was not obtained; references to its
question's wording concern the official hosted pre-publication version.

The companion archive supplies the source, frozen candidate, dated mathematical
and priority audits, bibliographic receipts, and reproducible finite and
integrity controls. Computation does not certify the infinite stochastic proof
or exhaustive priority. AI tools were used extensively in development, writing,
reproduction, verification and adversarial review. This is an unrefereed
research note: it has not undergone external human peer review or formal
proof-assistant certification. The note and companion materials are released
under CC BY 4.0.

\begin{thebibliography}{9}\small\raggedright
\bibitem{CH} Y. Couzini\'e and C. Hirsch,
\emph{Weakly reinforced P\'olya urns on countable networks},
Electronic Communications in Probability \textbf{26} (2021), paper 35, 1--10.
\href{https://doi.org/10.1214/21-ECP404}{doi:10.1214/21-ECP404}.
Inspected author version: \href{https://arxiv.org/abs/2010.03347v2}{arXiv:2010.03347v2}.
\bibitem{thesis} Y. Couzini\'e,
\emph{Sublinearly reinforced P\'olya urns on graphs of bounded degree},
master's thesis, Ludwig-Maximilians-Universit\"at M\"unchen, 17 September 2018.
\href{https://www.theorie.physik.uni-muenchen.de/TMP/theses/couziniethesis.pdf}{Primary thesis}.
\bibitem{OWR} V. Kleptsyn (joint work with C. Hirsch and M. Holmes),
\emph{Graph-based interacting P\'olya urns}, in MATRIX-MFO Tandem Workshop:
Stochastic Reinforcement Processes and Graphs, Oberwolfach Reports
\textbf{20} (2023), 653--656, especially 654--655.
\href{https://doi.org/10.4171/OWR/2023/12}{Report doi:10.4171/OWR/2023/12}.
\bibitem{MATRIX} G. Aletti et al.,
\emph{Open problems in stochastic reinforcement models},
2023 MATRIX Annals, Springer (2025), 477--498.
\href{https://doi.org/10.1007/978-3-031-76738-8_29}{doi:10.1007/978-3-031-76738-8\_29}.
Inspected \href{https://www.matrix-inst.org.au/wp_Matrix2016/wp-content/uploads/2023/Aletti.pdf}{official pre-publication version}, \S5, 8--10.
\bibitem{WARM} R. van der Hofstad, M. Holmes, A. Kuznetsov and W. Ruszel,
\emph{Strongly reinforced P\'olya urns with graph-based competition},
Annals of Applied Probability \textbf{26} (2016), 2494--2539.
Inspected \href{https://arxiv.org/abs/1406.0449v1}{arXiv:1406.0449v1}, \S3.6, 28.
\bibitem{Qin} S. Qin,
\emph{Homogenization of weakly reinforced P\'olya urns on countable networks},
\href{https://arxiv.org/abs/2610.11780v1}{arXiv:2610.11780v1}, submitted 8 October 2026.
\bibitem{candidate} A. Kriebel,
\emph{Subcritical reinforcement: audited corrected positive-equilibrium candidate},
\href{https://github.com/AlecKriebel/Math/pull/764}{Math repository PR 764}.
Current frozen head: \texttt{88acf20cdd5b4fe8605dab7eb27733b7deda1c47}.
\end{thebibliography}
\end{document}
